\section{Intrinsic geometry of the contraction}\label{sec:intrinsic}

\subsection{Excluding paths outside the chart}

Choose the radius $R$ inside a larger common chart and small enough
that~\eqref{eq:main-tensor} compares lengths by factors between
$1/2$ and $2$. The closed model set $\{r_\tau\le R\}$ is compact:
it is the compact quotient of a closed bounded finite-dimensional
moment level. Its image in the Hausdorff unitary-gauge moduli space
is compact and closed. A continuous path leaving the open radius-$R$
image must therefore first meet its radius-$R$ boundary.

Until this first exit, Lemma~\ref{lem:short-paths} and the tensor
estimate give
\begin{equation}\label{eq:radial-cost}
|r_\tau(q)-r_\tau(p)|\le2\,\operatorname{length}_{g_t}(\gamma).
\end{equation}
These inequalities extend through finite-length lower-stratum
endpoints by continuity.

\begin{proof}[Proof of \eqref{eq:main-distance}]
Let both endpoints have radius at most $\rho$. Their model
connecting path from Lemma~\ref{lem:short-paths} has length at most
$L\rho$ and stays inside that radius. It gives an actual competitor
of length at most $2L\rho$. An escaping path pays at least
$(R-\rho)/2$ by~\eqref{eq:radial-cost}. For $\rho\le cR$ with
$c$ sufficiently small, such a path cannot improve the infimum.
Every remaining actual near-minimizer can be taken of length at
most $4L\rho$ and hence stays at radius at most $(1+8L)\rho$.

Conversely, a global model near-minimizer can be chosen with length
at most $2L\rho$. The model radius inequality keeps it within
$(1+2L)\rho$. On these enlarged but small domains the sharp
relative tensor error is $O(|t|+\rho^2)$. Apply it to both sets
of paths and take infima, obtaining~\eqref{eq:main-distance}.
No minimizer need exist: use paths within an arbitrary positive
tolerance and pass to zero tolerance.

For arbitrarily close endpoints, the coarse radius bound first
confines all candidates; the sharp multiplicative comparison is
then applied to their entire lengths. Thus there is no additive
$O(\rho)$ error and no division by their possibly tiny distance.
The model distance separates points, so the actual distance also
separates points in this neighborhood.
\end{proof}

\subsection{Distances to the exceptional loci}

\begin{proof}[Proof of Theorem~\ref{thm:distances}]
The model paths of Proposition~\ref{prop:model-distances} decrease
radius and have length at most the starting radius $r$. Applying
\eqref{eq:main-tensor} gives the upper bounds in
\eqref{eq:distance-zero}--\eqref{eq:distance-rank}; their endpoints
lie in the stated local target sets.

Any actual competitor capable of lowering these bounds can be
taken of length at most $2r$. For $r\le cR$ it cannot leave the
chart, by the same exit-cost argument, and
\eqref{eq:radial-cost} confines it to radius $5r$. Its inverse
image ends in the appropriate model rank locus, so its model
length is bounded below by the \emph{global} distance in
\eqref{eq:model-distances}. The sharp tensor comparison on
radius $5r$ proves the lower bounds. This remains multiplicative
when the smaller distance is arbitrarily small compared with $r$.

On the zero section, $r_\tau^2=2|\tau|$, so the induced metric
comparison has relative error $O(|t|)$. On the rank-one fibre
over singular value $s$, equation~\eqref{eq:radius} gives
$r_\tau^2=\sqrt{\tau^2+4s^2}+|\tau|=O(|t|+s)$.
The exact induced fibre metrics computed in
Section~\ref{sec:model} therefore give the asserted
Grassmannian and projective-plane formulas.
\end{proof}

The collapse length scale is consequently
$\sqrt{|\tau|}=\sqrt{28\pi|t|/3}$. Both the deepest Grassmannian
and the $\PP^2$ fibres over rank one have this scale, even while
the rank-one point approaches the vertex at an unrelated rate.
The error in the actual metric remains controlled by the total
deformation radius, without requiring a positive lower bound for
the smaller transverse distance.

\subsection{The complete inner metric and outer overlap}

\begin{proof}[Proof of Corollary~\ref{cor:limits}]
Put $\kappa=|\tau|$ and write $x=\sqrt\kappa\,\widehat x$.
The model homothety gives
$g_\tau^{\rm q}=\kappa g_{\pm1}^{\rm q}$ and
$r_\tau^2=\kappa\widehat r^2$. Since
$|t|=3\kappa/(28\pi)$, substitution into
\eqref{eq:main-tensor} proves~\eqref{eq:inner-intro}.
Each bounded-radius region of the complete model is eventually
included; the same holds on any expanding radius $R_t$ with
$\kappa R_t^2\to0$.

For the pointed metric statement, fix a zero-section base point.
The physical cost of leaving the radius-$R$ chart stays bounded
below independently of $t$. After division of lengths by
$\sqrt\kappa$, that cost diverges. Equation~\eqref{eq:radial-cost}
also bounds $\widehat r$ on every bounded rescaled actual ball
about the base point. Thus all such balls lie in controlled
compact subcharts. On them~\eqref{eq:main-distance} has relative
error $O(\kappa)$, and every fixed bounded model ball is covered.
These estimates give pointed Gromov--Hausdorff convergence of
bounded balls. If balls of exactly equal radius are used, the
$1+O(\kappa)$ distance distortion changes radius only by
$O(\kappa)$; move along a model minimizing geodesic to obtain
the usual correspondence. Properness and completeness of the
model were proved in Proposition~\ref{prop:finite-geometry}.
No completeness of the full actual bundle-moduli component
is required.

On the common rank-two complex locus, Lemma~\ref{lem:outer-model}
gives relative error $O(|\tau|/s_2)$. If that ratio is bounded
and small, then~\eqref{eq:radius} gives $r_\tau^2=O(r_0^2)$.
Composition with~\eqref{eq:main-tensor} proves
\eqref{eq:outer-intro} first relative to $g_0^{\rm q}$.
Applying~\eqref{eq:main-tensor} also at $t=0$ compares the
actual chamber and actual wall metrics in their common $N$
coordinates. The overlap region
$|\tau|\ll s_2\le s_1\ll1$ matches the complete inner model
to the wall. At $s_2$ comparable to or smaller than $|\tau|$,
the inner estimate and full distance profile apply instead.
\end{proof}

All convergence and matching assertions above are statements about
quadratic forms, lengths and metric spaces. Differentiated estimates
would be needed for curvature or smooth convergence. The argument
also uses a fixed smooth bundle near its reducible splitting; it
does not include a non-locally-free or bubbling boundary. Within
this neighborhood, however, the entire contraction and its change
of stabilizer are described by the actual harmonic metric with
the same relative error.
